% !TeX root=main.tex
% در این فایل، عنوان پایان‌نامه، مشخصات خود و چکیده پایان‌نامه را به انگلیسی، وارد کنید.

%%%%%%%%%%%%%%%%%%%%%%%%%%%%%%%%%%%%
\baselineskip=.6cm
\begin{latin}
\latinuniversity{Iran University of Science and Technology}
\latinfaculty{School of Computer Engineering}
\latinsubject{Computer Engineering}
\latinfield{Computer Networks}
\latintitle{Collaborative Development of a Mobile Application to Support Speech Development in Children Under Three Years of Age}
\firstlatinsupervisor{Dr. Zeinab Movahedi}
%\secondlatinsupervisor{Second Supervisor}
% \firstlatinadvisor{First Advisor}
%\secondlatinadvisor{Second Advisor}
\latinname{Mohammadreza}
\latinsurname{Saemipour}
\latinthesisdate{June 2026}
\latinkeywords{Child Speech Development, Vocabulary Learning, Audiovisual Learning, Gamification, Mobile Application, Kotlin Multiplatform}
\en-abstract{
	The first three years of life are a critical period for developing speech skills and expanding a child's vocabulary. Continuous language interaction suited to the child's abilities can support this process; however, parents do not always have access to structured, engaging, and customizable content for daily practice. This project contributes to the design and development of NoVazheh, a mobile application intended to support vocabulary learning and create interactive opportunities between parents and children under three years of age. Persian words are organized into thematic categories and presented through images and audio. After hearing a word, the child selects the corresponding image and receives visual and auditory feedback. Questions prioritize words that have not yet been learned, while interactive games aim to sustain motivation and strengthen attention, memory, and visual discrimination. The parent section supports child profile management, custom categories and words using selected images and recorded audio, usage-time limits, and child mode. Children's responses are recorded to calculate the number of attempts, correct responses, overall accuracy, and progress in each category. NoVazheh uses a client-server architecture: a shared Android and iOS user interface is implemented with Kotlin and Compose Multiplatform, while a Ktor server manages content, authentication, and learning progress. The result is a functional, extensible multiplatform prototype that integrates audiovisual learning, gamification, content personalization, and progress tracking in a child-centered environment. It is intended as a supplementary tool alongside direct parent-child interaction.}
\latinfirstPage
\end{latin}
